\documentclass[11pt,a4paper]{article}
\usepackage{notes}

\title{Geophysics 3A --- Week 8 Workshop}
\author{Geothermics: Steady-State Heat Flow}
\date{\today}

\begin{document}
\maketitle

\section*{Workshop Objectives}
By the end of this workshop, students should be able to:
\begin{itemize}
  \item Use conservation of energy to reason about steady-state geotherms.
  \item Explain how thermal conductivity contrasts redistribute isotherms.
  \item Interpret geotherm curvature in terms of crustal heat production.
  \item Relate conductivity and heat production variations to geothermal potential.
\end{itemize}

\section*{0:00--0:15 \quad Conceptual Framing}
\begin{itemize}
  \item Steady-state heat flow as an energy bookkeeping problem.
  \item Review Fourier's law and the meaning of heat flow.
  \item Emphasize the separation of roles:
  \begin{itemize}
    \item conductivity controls \emph{gradients},
    \item heat production controls \emph{curvature}.
  \end{itemize}
\end{itemize}

\section*{0:15--0:45 \quad Refraction Exercise I: Insulating Sediments}
\paragraph{Setup}
Compare two crustal columns with identical basal heat flow:
\begin{itemize}
  \item Case A: crystalline crust exposed at the surface.
  \item Case B: crystalline crust overlain by a low-conductivity sedimentary layer.
\end{itemize}

\paragraph{Activities}
\begin{itemize}
  \item Predict temperature gradients in sediments vs basement.
  \item Sketch qualitative isotherms for each case.
  \item Identify where isotherms compress and where they stretch.
\end{itemize}

\paragraph{Key Concepts}
\begin{itemize}
  \item Heat flow conserved across layers.
  \item Lower conductivity $\rightarrow$ steeper gradient.
  \item Sediments act as a thermal blanket.
\end{itemize}

\paragraph{Geological Discussion}
\begin{itemize}
  \item Implications for geothermal gradients measured in boreholes.
  \item Relative geothermal potential of sedimentary basins vs exposed cratons.
\end{itemize}

\section*{0:45--1:20 \quad Refraction Exercise II: Layered Crustal Materials}
\paragraph{Focus}
Thermal refraction arising from conductivity contrasts within the crust.

\paragraph{Examples}
\begin{itemize}
  \item Shales (low $k$) vs quartzites (high $k$).
  \item Cyclic sedimentary layering preserved through metamorphism.
\end{itemize}

\paragraph{Activity}
\begin{itemize}
  \item Apply $q = -k \nabla T$ across layered media.
  \item Examine how gradients vary layer by layer.
  \item Discuss effective anisotropy in cyclically layered strata.
\end{itemize}

\paragraph{Connection}
\begin{itemize}
  \item Tie back to Week 2 physical-properties mixing demos.
  \item Discuss refraction in tilted and folded strata.
\end{itemize}

\section*{1:20--1:30 \quad Break}

\section*{1:30--2:05 \quad Energy Conservation and Geotherm Curvature}
\paragraph{Central Question}
How can two terranes with identical surface heat flow have different Moho temperatures?

\paragraph{Steps}
\begin{itemize}
  \item Write the energy balance:
  \[
    q_0 = q_b + \int_0^{z_c} A(z)\,dz
  \]
  \item Infer relative crustal heat production given different basal inputs.
  \item Introduce the steady-state solution and the curvature term:
  \[
    T(z) = T_0 + \frac{q_0}{k}z - \frac{A}{2k}z^2
  \]
\end{itemize}

\paragraph{Key Insight}
Curvature differences imply different deep temperatures even when surface gradients match.

\section*{2:05--2:40 \quad Why Is Archean Heat Production So Low?}
\paragraph{Observation}
Decay alone predicts high early heat production, yet preserved Archean crust shows limits.

\paragraph{Guided Discussion Topics}
\begin{itemize}
  \item Differentiation and partitioning of LILEs (K, Th, U).
  \item Selective erosion of felsic upper crust.
  \item Reworking: remelting and redistribution of older crust.
  \item Foundering of dense mafic roots.
  \item Secular mantle thermal evolution.
\end{itemize}

\section*{2:40--3:00 \quad Synthesis}
\begin{itemize}
  \item Summarize roles of $k$ and $A$.
  \item Connect observations to upcoming steady-state problem set.
  \item Preview transition to transient heat flow.
\end{itemize}

\end{document}